% Chapter 4, Section 2 _Linear Algebra_ Jim Hefferon
%  http://joshua.smcvt.edu/linearalgebra
%  2001-Jun-12
\section{相似}
\index{similarity|(}
%<*SimilarityMotiviation0>
我们已经定义：若存在非奇异的 \( P \) 与 \( Q \)，使得 \( \hat{H}=PHQ \)，
则称矩阵 \( H \) 与 \( \hat{H} \)
是矩阵等价的。
下图给出了这样做的动机，
图中 $H$ 与 $\hat{H}$ 都表示同一个映射
$h$，只是所相对于的
基对不同，分别为 $B,D$ 与 $\hat{B},\hat{D}$。
\begin{equation*}
  \begin{CD}
    V_{\wrt{B}}                   @>h>H>        W_{\wrt{D}}       \\
    @V{\scriptstyle\identity} VV             @V{\scriptstyle\identity} VV \\
    V_{\wrt{\hat{B}}}             @>h>\hat{H}>  W_{\wrt{\hat{D}}}
  \end{CD}
\end{equation*}
%</SimilarityMotiviation0>

%<*SimilarityMotiviation1>
现在考虑变换这一特殊情形，
此时陪域等于定义域，我们再补充要求
陪域的基等于定义域的基。
于是，我们所考虑的是相对于
$B,B$ 与 $D,D$ 的表示。\index{arrow diagram}
\begin{equation*}
  \begin{CD}
    V_{\wrt{B}}                   @>t>T>        V_{\wrt{B}}       \\
    @V{\scriptstyle\identity} VV              @V{\scriptstyle\identity} VV \\
    V_{\wrt{D}}                   @>t>\hat{T}>        V_{\wrt{D}}
  \end{CD}
\end{equation*}
用矩阵的语言来说，
\(\rep{t}{D,D}
  =\rep{\identity}{B,D}\;\rep{t}{B,B}\;\bigl(\rep{\identity}{B,D}\bigr)^{-1} \)。
%</SimilarityMotiviation1>



\subsection{定义与例题}
\begin{example}
考虑求导变换
$\map{d/dx}{\polyspace_2}{\polyspace_2}$，
以及该空间的两组基
$B=\sequence{x^2,x,1}$ 与
$D=\sequence{1,1+x,1+x^2}$。
我们将计算这个箭头方图的四条边。
\begin{equation*}
  \begin{CD}
    {\polyspace_2\,}_{\wrt{B}}                   @>d/dx>T>        {\polyspace_2\,}_{\wrt{B}}       \\
    @V{\scriptstyle\identity} VV              @V{\scriptstyle\identity} VV \\
    {\polyspace_2\,}_{\wrt{D}}                   @>d/dx>\hat{T}>        {\polyspace_2\,}_{\wrt{D}}
  \end{CD}
\end{equation*}
先看上边。
该变换作用于起始基~$B$ 的效果
\begin{equation*}
  x^2\mapsunder{d/dx} 2x
  \qquad
  x\mapsunder{d/dx} 1
  \qquad
  1\mapsunder{d/dx} 0
\end{equation*}
再相对于终止基（也是~$B$）表示
\begin{equation*}
  \rep{2x}{B}=\colvec{0 \\ 2 \\ 0}
  \qquad
  \rep{1}{B}=\colvec{0 \\ 0 \\ 1}
  \qquad
  \rep{0}{B}=\colvec{0 \\ 0 \\ 0}
\end{equation*}
便得到了该映射
的表示。
\begin{equation*}
  T=
  \rep{d/dx}{B,B}=
  \begin{mat}
    0 &0 &0 \\
    2 &0 &0 \\
    0 &1 &0
  \end{mat}
\end{equation*}

接下来，计算右边的矩阵需要
求出恒等映射作用于~$B$ 中元素的效果。
当然，恒等映射对它们毫无改变，因此
为得到该矩阵，我们把 $B$ 的元素相对于
~$D$ 表示。
\begin{equation*}
  \rep{x^2}{D}=\colvec{-1 \\ 0 \\ 1}
  \quad
  \rep{x}{D}=\colvec{-1 \\ 1 \\ 0}
  \quad
  \rep{1}{D}=\colvec{1 \\ 0 \\ 0}
\end{equation*}
于是，沿右边向下的矩阵就是将这些列向量依次并置而成的。
\begin{equation*}
  P=\rep{id}{B,D}=
  \begin{mat}
    -1 &-1  &1 \\
     0 &1   &0 \\
     1 &0   &0   
  \end{mat}
\end{equation*}

有了它，我们有两种方法来计算沿左边向上
的矩阵。
直接计算是把~$D$ 的元素相对于
~$B$ 表示，
\begin{equation*}
  \rep{1}{B}=\colvec{0 \\ 0 \\ 1}
  \quad
  \rep{1+x}{B}=\colvec{0 \\ 1 \\ 1}
  \quad
  \rep{1+x^2}{B}=\colvec{1 \\ 0 \\ 1}
\end{equation*}
再把它们并置成矩阵。
\begin{equation*}
  \begin{mat}
    0 &0 &1 \\
    0 &1 &0 \\ 
    1 &1 &1
  \end{mat}
\end{equation*}
计算沿左边向上的矩阵的另一种方法，
是取沿右边向下的矩阵~$P$ 的
逆矩阵。
\begin{equation*}
  \begin{pmat}{ccc|ccc}
    -1 &-1 &1 &1 &0 &0 \\
     0 &1  &0 &0 &1 &0 \\
     1 &0  &0 &0 &0 &1
  \end{pmat}
  % \grstep{\rho_1+\rho_3}
  % \grstep{\rho_2+\rho_3}
  % \grstep{-\rho_1}
  % \grstep{\rho_3+\rho_1}
  % \grstep{-\rho_2+\rho_1}
  \grstep{}
  \cdots
  \grstep{}
  \begin{pmat}{ccc|ccc}
     1 &0  &0 &0 &0 &1 \\
     0 &1  &0 &0 &1 &0 \\
     0 &0  &1 &1 &1 &1
  \end{pmat}
\end{equation*}

剩下方块的底边。
计算矩阵~$\hat{T}$ 有两种方法。
一种是直接计算它，
即求出该变换作用于~$D$ 中元素的效果，
\begin{equation*}
  1\mapsunder{d/dx} 0
  \qquad
  1+x\mapsunder{d/dx} 1
  \qquad
  1+x^2\mapsunder{d/dx} 2x
\end{equation*}
再相对于~$D$ 表示。
\begin{equation*}
  \hat{T}=
  \rep{d/dx}{D,D}=
  \begin{mat}
    0 &1 &-2 \\
    0 &0 &2 \\
    0 &0 &0
  \end{mat}
\end{equation*}
计算~$\hat{T}$ 的另一种方法——这也是我们通常采用的方法——
是沿着图先向上，再向左，然后向下。
\begin{align*}
  \rep{d/dx}{D,D}
  &=\rep{\identity}{B,D}\,\rep{d/dx}{B,B}\,\rep{\identity}{D,B}  \\
  \hat{T}
  &=\rep{\identity}{B,D}\,T\,\rep{\identity}{D,B}   \\
  &=  
  \begin{mat}
    -1 &-1  &1 \\
     0 &1   &0 \\
     1 &0   &0   
  \end{mat}
  \begin{mat}
    0 &0 &0 \\
    2 &0 &0 \\
    0 &1 &0
  \end{mat}
  \begin{mat}
    0 &0 &1 \\
    0 &1 &0 \\ 
    1 &1 &1
  \end{mat}
\end{align*}
展开乘积，得到的矩阵~$\hat{T}$ 与上面求出的相同。
\end{example}

\begin{definition} \label{df:Similar}
%<*df:Similar>
矩阵 \( T \) 与 $\hat{T}$ 是
\definend{相似的}\index{matrix!similarity}%
\index{equivalence relation!matrix similarity}\index{similar matrices}
若存在非奇异的 \( P \) 使得
$
  \hat{T}=PTP^{-1}
$。
%</df:Similar>
\end{definition}

\noindent 由于非奇异矩阵是方阵，
$T$ 与 $\hat{T}$ 必须
都是同尺寸的方阵。
\nearbyexercise{exer:SimIsEquivRel} 验证了
相似是一个等价关系。

\begin{example}
该定义并不要求我们考虑某个映射。
对这两个矩阵进行计算
\begin{equation*}
  P=
  \begin{mat}[r]
    2  &1  \\
    1  &1
  \end{mat}
  \qquad
  T=
  \begin{mat}[r]
    2  &-3  \\
    1  &-1
  \end{mat}
\end{equation*}
可知 $T$ 相似于这个矩阵。
\begin{equation*}
  \hat{T}=
  \begin{mat}[r]
    12  &-19  \\
    7  &-11
  \end{mat}
\end{equation*}
\end{example}

\begin{example}  \label{ex:OnlyZeroSimToZero}
%<*ex:OnlyZeroSimToZero>
与零矩阵相似的矩阵只有它自身：~$PZP^{-1}=PZ=Z$。
单位矩阵也有同样的性质：~$PIP^{-1}=PP^{-1}=I$。
%</ex:OnlyZeroSimToZero>
\end{example}

一个常见的特殊情形是向量空间为 $\C^n$，且
矩阵~$T$ 是某个映射相对于标准基的表示。
\begin{equation*}
  \begin{CD}
    \C^n_{\wrt{E_n}}                   @>t>T>        \C^n_{\wrt{E_n}}       \\
    @V{\scriptstyle\identity} VV              @V{\scriptstyle\identity} VV \\
    \C^n_{\wrt{D}}                   @>t>\hat{T}>        \C^n_{\wrt{D}}
  \end{CD}
\end{equation*}
此时在相似等式
$\hat{T}=PTP^{-1}$ 中，~$P$ 的列就是~$D$ 的元素。

矩阵相似是矩阵等价的特殊情形，因此
若两个矩阵相似，则它们矩阵等价。
反过来又如何呢：~若它们是方阵，
任意两个矩阵等价的矩阵一定相似吗？
不；一个单位矩阵的矩阵等价类
由该尺寸的所有非奇异矩阵组成，
而前面的例题表明一个单位矩阵的相似类
中唯一的成员是它自身。
因此这两个矩阵
矩阵等价但不相似。
\begin{equation*}
  T=
  \begin{mat}[r]
    1  &0  \\
    0  &1
  \end{mat}
  \qquad
  S=
  \begin{mat}[r]
    1  &2  \\
    0  &3
  \end{mat}
\end{equation*}
于是，某些矩阵等价类
会分裂成两个或更多个相似类\Dash 相似给出了比矩阵等价更细的
划分。
这里展示了一些矩阵等价类被细分为
相似类。
\begin{center}
  \includegraphics{jc/mp/ch5.4}
\end{center}

为了理解相似关系，我们将研究相似类。
我们处理这个问题的方式，与我们此前研究
行等价关系与矩阵等价关系的方式相同，即求出
相似类的代表元的一种
规范形式，% \appendrefs{representatives}\spacefactor=1000 %
称为 Jordan 形式。
有了这种规范形式，我们就能通过检查
两个矩阵是否处于以同一个代表元为代表的类中，来判定它们是否相似。
对于行等价与矩阵等价，我们都已看到，一种
规范形式能让我们洞见同一类中的成员
在哪些方面彼此相似
（例如，两个同尺寸的矩阵
当且仅当它们有相同的秩时矩阵等价）。
%(Along the way we shall see ideas that are interesting and important
%in their own right, not just as stepping stones to Jordan form.)

